Misled by superficial coincidences, [Jammer] imagines that Bohr's thought has been influenced, through Høffding, by Kierkegaard and William James. There can be no doubt that his surmise is unfounded. Bohr was a completely independent thinker; from early youth, he developed his epistemological ideas single-handed and with no more philosophical preparation than H0ffding's elemen- tary course of lectures. What the latter, and still more his father, encouraged in him was an attitude of open-mindedness and unprejudiced approach to the problems, and certainly not any reliance on particular philosophical systems. There was indeed one person who did exert a real influence on Bohr's views: this was his brother Harald. Niels Bohr told me once how much he owed to Harald for the clarification of his ideas about the function of mathematics, which were of great moment in the early formation of his more general conception of the role of language. If anything, Harald stronglydissuaded Niels from taking any interest in Kierkegaard. As to Williams James, it was only after the completion of Bohr's epistemological analysis of quantum mechanics that his attention was called to the Principles of Psychology: the similarity of attitudes interested Bohr, but did not influence him.

The author could have inferred all this (or most of it) from published statements; but he somehow went astray and ponderously built in a completely fictitious "Kierkegaard-Høffding" ideology into the discussion of Bohr's work. This unfortunate error runs through the whole book and spoils the argument wherever it intrudes into it. Moreover, he has failed to realize that Bohr's conception of complementarity in quantum mechanics is not the expression of a "specific philosophical position", but an inherent part of the theory which has the same validity as its formal aspect and is inseparable from it.




Ved betragtning af den resignation, til hvilken den beskrevne situation tvinger os med hensyn til de ønsker om anskuelighed, som giver vort hele sprog dets præg, er det meget lærerigt, at vi allerede ved simple, psykologiske erfaringer møder grundtræk ikke alene af den relativistiske, men også af den reciprokke betragtningsmåde. Til den relativitet af vore bevægelsesfænomener, hvormed enhver allerede fra barndommen er fortrolig gennem kørsel og sejlads, svarer dagligdags erfaringer over berøringsfornemmelsernes reciprocitet. Vi skal her blot minde om den af psykologer ofte fremdragne, ejendommelige fornemmelse, som enhver har oplevet ved forsøget på at orientere sig i et mørkt værelse ved at famle sig frem med en stok. Medens stokken, så længe vi holder let på den, fremtræder som genstand for berøringssansen, taber vi ved et fast tag om den indtrykket af et fremmed legeme, og berøringsfornemmelsen henføres umiddelbart til det punkt, hvor stokken støder på de legemer vi vil undersøge. Det er næppe nogen overdrivelse, når man allerede ud fra psykologiske erfaringer vil påstå, at begreberne tid og rum efter deres væsens art først får mening gennem muligheden for at se bort fra vekselvirkningen med målemidlerne. I det hele bringer analysen af sanseindtrykkene en bemærkelsesværdig uafhængighed for dagen, hvad det psykologiske grundlag for rum-og tidsforestillingerne på den ene side og de på kraftvirkninger begrundede forestillinger om energi og bevægelsesmængde på den anden side angår. Imidlertid kendetegnes dette område, som allerede berørt, fremfor alt ved reciprocitetsforhold, som hænger sammen med bevidsthedens enhedskarakter og som viser en påfaldende lighed med de fysiske følger af virkningskvantet. Det drejer sig her om velkendte ejendommeligheder ved følelses-og viljelivet, som ganske unddrager sig en fremstilling ved anskuelige billeder. I særdeleshed udviser den tilsyneladende modsætning mellem den associative tænknings kontinuerte fremadskriden og bevarelsen af personlighedens enhed en tankevækkende lighed med forholdet mellem den af superpositionsprincippet beherskede bølgebeskrivelse for materielle smådeles bevægelser og deres uforstyrrelige individualitet. Den uundgåelige påvirkning af de atomare fænomener ved deres iagttagelse svarer her til den velkendte ændring af de psykiske oplevelsers farve, som ledsager enhver fæsten af opmærksomheden ved en enkelt af deres forskellige sider.


Det erkendelsesproblem, som det drejer sig om, er i korthed følgende: en beskrivelse af vor tankevirksomhed forlanger på den ene side, at et objektivt givet indhold stilles over for et betragtende subjekt, medens på den anden side – således som det fremgår allerede af et sådant udsagn – ingen skarp adskillelse mellem objekt og subjekt kan opretholdes, da jo også det sidstnævnte begreb tilhører vort tankeindhold.


%%% Local Variables:
%%% mode: latex
%%% TeX-master: t
%%% End:
